\label{sec:appendix}

\section{Forward Transfer} \label{sec:app_results_fwt}


\cref{fig:results_fwt} in the main paper shows the forward transfer plots for some datasets but not all due to space limitations. \cref{fig:app_results_fwt} presents the forward transfer results for all datasets. 
It can be clearly seen that SPG has the best positive forward transfer and keeps or even grows it constantly in all cases. On the other hand, the other parameter isolation-based methods, PGN, PathNet, HAT, CAT, and WSN lose their ability for the forward transfer in later tasks for the dissimilar task (i.e., (a) to (e)).


\begin{figure}[h]

\centering
\begin{minipage}{.27\hsize}
\centering
\includegraphics[width=\hsize]{images/fwt_C-10_leg.pdf}
\subcaption{C-10}
\end{minipage}
%\hfill
\begin{minipage}{.27\hsize}
\centering
\includegraphics[width=\hsize]{images/fwt_C-20.pdf}
\subcaption{C-20}
\end{minipage}
%\hfill
\begin{minipage}{.27\hsize}
\centering
\includegraphics[width=\hsize]{images/fwt_T-10.pdf}
\subcaption{T-10}
\end{minipage}

\medskip

\begin{minipage}{.27\hsize}
\centering
\includegraphics[width=\hsize]{images/fwt_T-20.pdf}
\subcaption{T-20}
\end{minipage}
%\hfill
\begin{minipage}{.27\hsize}
\centering
\includegraphics[width=\hsize]{images/fwt_I-100.pdf}
\subcaption{I-100}
\end{minipage}
%\hfill
\begin{minipage}{.27\hsize}
\centering
\includegraphics[width=\hsize]{images/fwt_FC-10.pdf}
\subcaption{FC-10}
\end{minipage}

\medskip

\begin{minipage}{.27\hsize}
\centering
\includegraphics[width=\hsize]{images/fwt_FC-20.pdf}
\subcaption{FC-20}
\end{minipage}
%\hfill
\begin{minipage}{.27\hsize}
\centering
\includegraphics[width=\hsize]{images/fwt_FE-10.pdf}
\subcaption{FE-10}
\end{minipage}
%\hfill
\begin{minipage}{.27\hsize}
\centering
\includegraphics[width=\hsize]{images/fwt_FE-20.pdf}
\subcaption{FE-20}
\end{minipage}

\caption{Forward transfer results. (a) to (e) are for dissimilar tasks, and (f) to (i) are for similar tasks.}
\label{fig:app_results_fwt}

\end{figure}








\section{Capacity Consumption at Each Layer} \label{sec:app_capacity_layerwise}

We show the percentage of parameters in the whole network that are fully blocked in \cref{tab:blocking_ratio} in the main text of the paper. Here \cref{tab:app_blocking_ratio_layerwise} presents the same result for each layer. We use AlexNet as the backbone, and it has three convolution layers followed by two fully-connected layers.

\begin{table}[p]
\centering
\caption{
The percentage of parameters that are completely blocked for each layer.
$T$ is the total number of tasks (e.g., $T = 10$ for C-10).
}
\label{tab:app_blocking_ratio_layerwise}

\vspace{2mm}

\begin{minipage}{.85\hsize}
\centering
\subcaption{Results for the 1st convolution layer.}
\resizebox{\columnwidth}{!}{
\begin{tabular}{ccrrrrrrrrr}
\toprule
& & \multicolumn{5}{c}{Dissimilar tasks} & \multicolumn{4}{c}{Similar tasks} \\ \cmidrule(lr){3-7} \cmidrule(lr){8-11}
$t$ & Model &
\multicolumn{1}{c}{C-10} & \multicolumn{1}{c}{C-20} & \multicolumn{1}{c}{T-10} & \multicolumn{1}{c}{T-20} & \multicolumn{1}{c}{I-100} &
\multicolumn{1}{c}{FC-10} & \multicolumn{1}{c}{FC-20} & \multicolumn{1}{c}{FE-10} & \multicolumn{1}{c}{FE-20}  \\
\midrule

\multirow{2}{*}{$1$}
& HAT & 
$25.0\std{5.3}$ & $24.4\std{8.4}$ & 
$31.6\std{6.0}$ & $14.4\std{5.5}$ &
$60.0\std{4.4}$ & 
$17.8\std{4.4}$ & $40.0\std{8.7}$ &
$75.0\std{3.1}$ & $67.2\std{9.9}$ \\
& \textbf{SPG} & 
$0.0\std{0.0}$ & $0.0\std{0.0}$ &
$0.0\std{0.0}$ & $0.0\std{0.0}$ &
$0.0\std{0.0}$ &
$0.0\std{0.0}$ & $0.0\std{0.0}$ &
$0.0\std{0.0}$ & $0.0\std{0.0}$ \\

\hdashline

\multirow{2}{*}{$T/2$}
& HAT & 
$77.8\std{2.7}$ & $98.4\std{1.7}$ &
$83.4\std{3.6}$ & $79.1\std{2.5}$ &
$100.0\std{0.0}$ &
$73.4\std{7.6}$ & $100.0\std{0.0}$ &
$98.8\std{1.2}$ & $99.4\std{0.8}$ \\
& \textbf{SPG} & 
$0.0\std{0.0}$ & $0.0\std{0.0}$ &
$0.0\std{0.0}$ & $0.0\std{0.0}$ &
$0.0\std{0.0}$ &
$0.0\std{0.0}$ & $0.0\std{0.0}$ &
$0.0\std{0.0}$ & $0.0\std{0.0}$ \\

\hdashline

\multirow{2}{*}{$T$}
& HAT & 
$95.6\std{1.5}$ & $100.0\std{0.0}$ &
$97.2\std{1.2}$ & $95.3\std{1.0}$ &
$100.0\std{0.0}$ & 
$98.8\std{1.8}$ & $100.0\std{0.0}$ &
$100.0\std{0.0}$ & $100.0\std{0.0}$ \\
& \textbf{SPG} & 
$0.0\std{0.0}$ & $0.0\std{0.0}$ &
$0.0\std{0.0}$ & $0.0\std{0.0}$ &
$0.0\std{0.0}$ &
$0.0\std{0.0}$ & $0.0\std{0.0}$ &
$0.0\std{0.0}$ & $0.0\std{0.0}$ \\

\bottomrule
\end{tabular}
}
\end{minipage}

\vspace{2mm}

\begin{minipage}{.85\hsize}
\centering
\subcaption{Results for the 2nd convolution layer.}
\resizebox{\columnwidth}{!}{
\begin{tabular}{ccrrrrrrrrr}
\toprule
& & \multicolumn{5}{c}{Dissimilar tasks} & \multicolumn{4}{c}{Similar tasks} \\ \cmidrule(lr){3-7} \cmidrule(lr){8-11}
$t$ & Model &
\multicolumn{1}{c}{C-10} & \multicolumn{1}{c}{C-20} & \multicolumn{1}{c}{T-10} & \multicolumn{1}{c}{T-20} & \multicolumn{1}{c}{I-100} &
\multicolumn{1}{c}{FC-10} & \multicolumn{1}{c}{FC-20} & \multicolumn{1}{c}{FE-10} & \multicolumn{1}{c}{FE-20}  \\
\midrule

\multirow{2}{*}{$1$}
& HAT & 
$5.6\std{1.6}$ & $7.3\std{2.9}$ & 
$8.9\std{2.3}$ & $1.5\std{1.1}$ &
$33.9\std{5.0}$ & 
$2.2\std{0.6}$ & $12.0\std{2.0}$ &
$55.7\std{6.6}$ & $45.7\std{12.1}$ \\
& \textbf{SPG} & 
$0.1\std{0.0}$ & $0.1\std{0.1}$ &
$0.0\std{0.0}$ & $0.1\std{0.0}$ &
$0.0\std{0.0}$ &
$0.0\std{0.0}$ & $0.0\std{0.0}$ &
$0.0\std{0.1}$ & $0.0\std{0.0}$ \\

\hdashline

\multirow{2}{*}{$T/2$}
& HAT & 
$52.4\std{2.5}$ & $98.0\std{1.7}$ &
$63.7\std{6.2}$ & $56.2\std{4.5}$ &
$100.0\std{0.0}$ &
$41.3\std{4.7}$ & $98.3\std{1.2}$ &
$98.4\std{1.0}$ & $99.2\std{1.0}$ \\
& \textbf{SPG} & 
$0.6\std{0.1}$ & $2.1\std{0.4}$ &
$0.3\std{0.1}$ & $0.9\std{0.3}$ &
$1.5\std{0.3}$ &
$0.1\std{0.1}$ & $0.2\std{0.1}$ &
$0.5\std{0.5}$ & $0.6\std{0.2}$ \\

\hdashline

\multirow{2}{*}{$T$}
& HAT & 
$81.0\std{2.0}$ & $99.8\std{0.3}$ &
$90.2\std{2.8}$ & $86.4\std{4.5}$ &
$100.0\std{0.0}$ & 
$86.9\std{2.2}$ & $99.4\std{0.3}$ &
$100.0\std{0.0}$ & $100.0\std{0.0}$ \\
& \textbf{SPG} & 
$1.1\std{0.2}$ & $3.1\std{0.5}$ &
$0.4\std{0.1}$ & $1.5\std{0.3}$ &
$1.8\std{0.4}$ &
$0.2\std{0.1}$ & $0.3\std{0.2}$ &
$0.8\std{0.7}$ & $0.9\std{0.3}$ \\

\bottomrule
\end{tabular}
}
\end{minipage}


\vspace{2mm}


\begin{minipage}{.85\hsize}
\centering
\subcaption{Results for the 3rd convolution layer.}
\resizebox{\columnwidth}{!}{
\begin{tabular}{ccrrrrrrrrr}
\toprule
& & \multicolumn{5}{c}{Dissimilar tasks} & \multicolumn{4}{c}{Similar tasks} \\ \cmidrule(lr){3-7} \cmidrule(lr){8-11}
$t$ & Model &
\multicolumn{1}{c}{C-10} & \multicolumn{1}{c}{C-20} & \multicolumn{1}{c}{T-10} & \multicolumn{1}{c}{T-20} & \multicolumn{1}{c}{I-100} &
\multicolumn{1}{c}{FC-10} & \multicolumn{1}{c}{FC-20} & \multicolumn{1}{c}{FE-10} & \multicolumn{1}{c}{FE-20}  \\
\midrule

\multirow{2}{*}{$1$}
& HAT & 
$4.8\std{1.3}$ & $8.7\std{1.7}$ & 
$7.5\std{1.1}$ & $0.7\std{0.3}$ &
$30.6\std{2.8}$ & 
$1.0\std{0.2}$ & $8.0\std{1.5}$ &
$43.7\std{6.5}$ & $36.4\std{6.2}$ \\
& \textbf{SPG} & 
$0.1\std{0.0}$ & $0.1\std{0.0}$ &
$0.1\std{0.0}$ & $0.1\std{0.0}$ &
$0.1\std{0.0}$ &
$0.1\std{0.0}$ & $0.2\std{0.1}$ &
$0.0\std{0.0}$ & $0.0\std{0.0}$ \\

\hdashline

\multirow{2}{*}{$T/2$}
& HAT & 
$42.6\std{3.6}$ & $98.7\std{0.7}$ &
$60.3\std{4.1}$ & $47.1\std{2.3}$ &
$100.0\std{0.0}$ &
$14.7\std{1.2}$ & $88.7\std{4.2}$ &
$98.0\std{0.7}$ & $99.1\std{0.4}$ \\
& \textbf{SPG} & 
$0.7\std{0.1}$ & $2.1\std{0.4}$ &
$0.3\std{0.1}$ & $1.0\std{0.2}$ &
$1.2\std{0.2}$ &
$0.5\std{0.1}$ & $1.0\std{0.3}$ &
$0.3\std{0.2}$ & $0.3\std{0.1}$ \\

\hdashline

\multirow{2}{*}{$T$}
& HAT & 
$66.8\std{2.7}$ & $99.7\std{0.5}$ &
$87.1\std{2.2}$ & $74.1\std{3.4}$ &
$100.0\std{0.0}$ & 
$43.1\std{4.0}$ & $95.3\std{1.9}$ &
$99.8\std{0.5}$ & $99.8\std{0.2}$ \\
& \textbf{SPG} & 
$1.2\std{0.2}$ & $3.1\std{0.7}$ &
$0.4\std{0.1}$ & $1.5\std{0.3}$ &
$1.8\std{0.2}$ &
$0.9\std{0.2}$ & $1.4\std{0.4}$ &
$0.5\std{0.3}$ & $0.5\std{0.1}$ \\

\bottomrule
\end{tabular}
}
\end{minipage}


\vspace{2mm}


\begin{minipage}{.85\hsize}
\centering
\subcaption{Results for the 4th fully-connected layer.}
\resizebox{\columnwidth}{!}{
\begin{tabular}{ccrrrrrrrrr}
\toprule
& & \multicolumn{5}{c}{Dissimilar tasks} & \multicolumn{4}{c}{Similar tasks} \\ \cmidrule(lr){3-7} \cmidrule(lr){8-11}
$t$ & Model &
\multicolumn{1}{c}{C-10} & \multicolumn{1}{c}{C-20} & \multicolumn{1}{c}{T-10} & \multicolumn{1}{c}{T-20} & \multicolumn{1}{c}{I-100} &
\multicolumn{1}{c}{FC-10} & \multicolumn{1}{c}{FC-20} & \multicolumn{1}{c}{FE-10} & \multicolumn{1}{c}{FE-20}  \\
\midrule

\multirow{2}{*}{$1$}
& HAT & 
$2.7\std{0.8}$ & $12.6\std{1.6}$ & 
$4.0\std{0.3}$ & $0.2\std{0.1}$ &
$25.7\std{1.4}$ & 
$0.3\std{0.0}$ & $6.5\std{0.9}$ &
$27.7\std{2.3}$ & $22.7\std{2.1}$ \\
& \textbf{SPG} & 
$0.1\std{0.0}$ & $0.1\std{0.0}$ &
$0.1\std{0.0}$ & $0.1\std{0.1}$ &
$0.0\std{0.0}$ &
$0.1\std{0.0}$ & $0.2\std{0.0}$ &
$0.1\std{0.0}$ & $0.1\std{0.0}$ \\

\hdashline

\multirow{2}{*}{$T/2$}
& HAT & 
$29.7\std{1.4}$ & $98.3\std{0.6}$ &
$47.7\std{2.1}$ & $29.4\std{2.2}$ &
$99.9\std{0.1}$ &
$3.8\std{0.4}$ & $70.9\std{3.7}$ &
$90.2\std{1.1}$ & $95.6\std{1.4}$ \\
& \textbf{SPG} & 
$0.8\std{0.2}$ & $2.4\std{0.2}$ &
$0.4\std{0.1}$ & $1.4\std{0.2}$ &
$1.6\std{0.2}$ &
$0.8\std{0.1}$ & $1.4\std{0.1}$ &
$0.6\std{0.2}$ & $0.8\std{0.1}$ \\

\hdashline

\multirow{2}{*}{$T$}
& HAT & 
$52.4\std{1.1}$ & $99.7\std{0.3}$ &
$72.3\std{1.6}$ & $50.6\std{2.5}$ &
$99.9\std{0.1}$ & 
$12.1\std{1.3}$ & $84.0\std{1.8}$ &
$98.5\std{0.4}$ & $98.2\std{0.7}$ \\
& \textbf{SPG} & 
$1.6\std{0.2}$ & $4.1\std{0.3}$ &
$0.7\std{0.1}$ & $2.5\std{0.2}$ &
$2.3\std{0.4}$ &
$1.5\std{0.2}$ & $2.2\std{0.3}$ &
$0.9\std{0.4}$ & $1.2\std{0.2}$ \\

\bottomrule
\end{tabular}
}
\end{minipage}


\vspace{2mm}


\begin{minipage}{.85\hsize}
\centering
\subcaption{Results for the 5th fully-connected layer.}
\resizebox{\columnwidth}{!}{
\begin{tabular}{ccrrrrrrrrr}
\toprule
& & \multicolumn{5}{c}{Dissimilar tasks} & \multicolumn{4}{c}{Similar tasks} \\ \cmidrule(lr){3-7} \cmidrule(lr){8-11}
$t$ & Model &
\multicolumn{1}{c}{C-10} & \multicolumn{1}{c}{C-20} & \multicolumn{1}{c}{T-10} & \multicolumn{1}{c}{T-20} & \multicolumn{1}{c}{I-100} &
\multicolumn{1}{c}{FC-10} & \multicolumn{1}{c}{FC-20} & \multicolumn{1}{c}{FE-10} & \multicolumn{1}{c}{FE-20}  \\
\midrule

\multirow{2}{*}{$1$}
& HAT & 
$1.3\std{0.4}$ & $17.9\std{1.0}$ & 
$2.1\std{0.2}$ & $0.1\std{0.0}$ &
$22.0\std{0.5}$ & 
$0.2\std{0.0}$ & $6.0\std{1.0}$ &
$20.8\std{0.5}$ & $16.2\std{1.3}$ \\
& \textbf{SPG} & 
$0.1\std{0.0}$ & $0.2\std{0.0}$ &
$0.1\std{0.0}$ & $0.0\std{0.0}$ &
$0.1\std{0.0}$ &
$0.1\std{0.0}$ & $0.1\std{0.0}$ &
$0.1\std{0.0}$ & $0.1\std{0.0}$ \\

\hdashline

\multirow{2}{*}{$T/2$}
& HAT & 
$17.6\std{1.5}$ & $98.2\std{0.1}$ &
$29.9\std{2.0}$ & $15.4\std{1.2}$ &
$99.8\std{0.0}$ &
$2.1\std{0.2}$ & $61.8\std{1.8}$ &
$84.6\std{0.8}$ & $92.8\std{1.8}$ \\
& \textbf{SPG} & 
$1.0\std{0.1}$ & $3.2\std{0.4}$ &
$0.5\std{0.1}$ & $1.5\std{0.2}$ &
$5.5\std{0.6}$ &
$0.7\std{0.0}$ & $1.2\std{0.2}$ &
$0.7\std{0.3}$ & $1.1\std{0.2}$ \\

\hdashline

\multirow{2}{*}{$T$}
& HAT & 
$34.9\std{2.0}$ & $99.6\std{0.1}$ &
$48.5\std{2.1}$ & $30.6\std{2.0}$ &
$99.9\std{0.1}$ & 
$6.1\std{1.1}$ & $76.4\std{0.6}$ &
$97.4\std{0.3}$ & $97.1\std{1.1}$ \\
& \textbf{SPG} & 
$2.2\std{0.3}$ & $5.8\std{0.7}$ &
$1.1\std{0.2}$ & $3.4\std{0.3}$ &
$8.6\std{0.7}$ &
$1.3\std{0.0}$ & $1.9\std{0.2}$ &
$1.2\std{0.6}$ & $1.8\std{0.3}$ \\

\bottomrule
\end{tabular}
}
\end{minipage}


\end{table}



As we described in \cref{sec:capacity}, SPG blocks much fewer parameters than what HAT does in all cases.
Additionally, we can see from \cref{tab:app_blocking_ratio_layerwise} the significant difference between HAT and SPG in their layer-wise tendency. 
HAT blocks more parameters in earlier layers (e.g., after learning task 5 of C-10, 77.8\% of parameters in the 1st convolution layer are completely blocked while 52.4\% of the ones in the 2nd convolution layer are), which is reasonable given that the earlier layers are supposed to extract basic features and thus changing their parameters without being blocked could easily cause more forgetting than in later layers. 
On the other hand, SPG contrarily tends to completely block more parameters in later layers (e.g., after learning task 5 of C-10, 0.0\% of parameters in the 1st convolution layer are completely blocked while 0.6\% of ones in 2nd convolution layer are). Since SPG computes parameters' importance based on their gradients with regard to the loss through normalization, this result implies that later layers are likely to have more parameters on which some of the tasks highly depend. It can be said that SPG keeps earlier layers alive with less blocking for better basic feature learning (i.e., leading to positive knowledge transfer) while it blocks some specific parameters in later layers that are supposed to be important for previous tasks, which is different from what HAT does. 





\newpage







\section{Quantitative Analysis on Cross-Head Importance (CHI)} \label{sec:app_chi}

We analyze in detail how CHI quantitatively contributes to suppressing parameter updates with the following four metrics.

\textbf{(1) Overwrite \underline{\textit{F}}requency at \underline{\textit{each}} task (F-each)}:
How often does the importance from CHI have a larger value than the one from the current task (1 means that it always happens)? It corresponds to cases where $\bm{\gamma}_i^{t,\tau} > \bm{\gamma}_i^{t,t}$ for any $\tau (1 \le \tau < t)$ in \cref{eq:gamma_i_t}. 
\\
\textbf{(2) Overwrite \underline{\textit{G}}ap at \underline{\textit{each}} task (G-each)}:
When the cases of F-each happen, how much is the difference of overwriting on average? It is defined by the average of $\max{(\bm{\gamma}_i^{t,1},\cdots,\bm{\gamma}_i^{t,t-1})} - \bm{\gamma}_i^{t,t}$.
\\
\textbf{(3) Overwrite \underline{\textit{F}}requency in \underline{\textit{total}} (F-total)}:
How often does the importance from CHI actually overwrite the accumulated importance through the maximum operation? It corresponds to cases where $\bm{\gamma}_i^{t,\tau} > \bm{\gamma}_i^{\le t-1}$ for any $\tau (1 \le \tau < t)$ in \cref{eq:gamma_i_le_t}. 
\\
\textbf{(4) Overwrite \underline{\textit{G}}ap in \underline{\textit{total}} (G-total)}:
When the cases of F-total happen, how much is the difference of overwriting on average? It is defined by the average of $\max{(\bm{\gamma}_i^{t,1},\cdots,\bm{\gamma}_i^{t,t-1})} - \bm{\gamma}_i^{\le t-1}$.

The result is presented in \cref{tab:app_ablation_chi}.
We can clearly observe that CHI adds more importance to some parameters (e.g., in C-10, about 15-42\% of parameters constantly update their accumulated importance by the ones from CHI), which is denoted by F-total. Since we introduce CHI to further mitigate forgetting by accumulating more importance, this expectation is consistent with the observed results. Although CHI also overwrites the accumulated importance in similar tasks as frequently as in dissimilar tasks (see F-each and F-total), it happens with a smaller gap overall (see G-each and G-total), which is reasonable as the tasks are similar thus parameters can have similar gradients among different tasks.


\begin{table}[tb]
\centering

\caption{
Quantitative contribution of CHI in learning task $t$. $T$ is the total number of tasks (e.g., $T=10$ for C-10). 
}
\label{tab:app_ablation_chi}

\resizebox{\columnwidth}{!}{
\begin{tabular}{crrrrrrrrrrrr}
\toprule
& \multicolumn{4}{c}{C-10} & \multicolumn{4}{c}{C-20} & \multicolumn{4}{c}{T-10} \\ 
\cmidrule(lr){2-5} \cmidrule(lr){6-9} \cmidrule(lr){10-13}
$t$ & 
\multicolumn{1}{c}{F-each} & \multicolumn{1}{c}{G-each} & \multicolumn{1}{c}{F-total} & \multicolumn{1}{c}{G-total} & 
\multicolumn{1}{c}{F-each} & \multicolumn{1}{c}{G-each} & \multicolumn{1}{c}{F-total} & \multicolumn{1}{c}{G-total} &
\multicolumn{1}{c}{F-each} & \multicolumn{1}{c}{G-each} & \multicolumn{1}{c}{F-total} & \multicolumn{1}{c}{G-total} \\
\midrule

$2$ &
% C-10
$0.64\std{0.02}$ & $0.15\std{0.02}$ & $0.42\std{0.02}$ & $0.13\std{0.00}$ &
% C-20
$0.52\std{0.09}$ & $0.10\std{0.03}$ & $0.35\std{0.06}$ & $0.09\std{0.02}$ &
% T-10
$0.49\std{0.01}$ & $0.13\std{0.00}$ & $0.32\std{0.01}$ & $0.12\std{0.00}$ \\

$T/2$ &
% C-10
$0.84\std{0.01}$ & $0.31\std{0.02}$ & $0.24\std{0.01}$ & $0.04\std{0.00}$ &
% C-20
$0.88\std{0.01}$ & $0.33\std{0.02}$ & $0.15\std{0.03}$ & $0.01\std{0.00}$ &
% T-10
$0.77\std{0.01}$ & $0.28\std{0.01}$ & $0.24\std{0.01}$ & $0.04\std{0.00}$ \\

$T$ &
% C-10
$0.90\std{0.03}$ & $0.39\std{0.04}$ & $0.15\std{0.01}$ & $0.01\std{0.00}$ &
% C-20
$0.90\std{0.00}$ & $0.34\std{0.01}$ & $0.05\std{0.00}$ & $0.00\std{0.00}$ &
% T-10
$0.87\std{0.01}$ & $0.36\std{0.01}$ & $0.13\std{0.00}$ & $0.01\std{0.00}$ \\

\bottomrule
\end{tabular}
}


\resizebox{\columnwidth}{!}{
\begin{tabular}{crrrrrrrrrrrr}
\toprule
& \multicolumn{4}{c}{T-20} & \multicolumn{4}{c}{I-100} & \multicolumn{4}{c}{FC-10} \\ 
\cmidrule(lr){2-5} \cmidrule(lr){6-9} \cmidrule(lr){10-13}
$t$ & 
\multicolumn{1}{c}{F-each} & \multicolumn{1}{c}{G-each} & \multicolumn{1}{c}{F-total} & \multicolumn{1}{c}{G-total} & 
\multicolumn{1}{c}{F-each} & \multicolumn{1}{c}{G-each} & \multicolumn{1}{c}{F-total} & \multicolumn{1}{c}{G-total} &
\multicolumn{1}{c}{F-each} & \multicolumn{1}{c}{G-each} & \multicolumn{1}{c}{F-total} & \multicolumn{1}{c}{G-total} \\
\midrule

$2$ &
% T-20
$0.52\std{0.01}$ & $0.14\std{0.01}$ & $0.34\std{0.02}$ & $0.12\std{0.01}$ &
% I-100
$0.47\std{0.01}$ & $0.11\std{0.00}$ & $0.30\std{0.02}$ & $0.11\std{0.01}$ &
% FC-10
$0.53\std{0.04}$ & $0.14\std{0.01}$ & $0.38\std{0.03}$ & $0.12\std{0.00}$ \\

$T/2$ &
% T-20
$0.88\std{0.02}$ & $0.37\std{0.02}$ & $0.14\std{0.00}$ & $0.01\std{0.00}$ &
% I-100
$0.91\std{0.00}$ & $0.36\std{0.00}$ & $0.02\std{0.00}$ & $0.00\std{0.00}$ &
% FC-10
$0.83\std{0.06}$ & $0.25\std{0.01}$ & $0.25\std{0.01}$ & $0.04\std{0.00}$ \\

$T$ &
% T-20
$0.92\std{0.01}$ & $0.41\std{0.02}$ & $0.05\std{0.00}$ & $0.00\std{0.00}$ &
% I-100
$0.91\std{0.01}$ & $0.37\std{0.00}$ & $0.02\std{0.00}$ & $0.00\std{0.00}$ &
% FC-10
$0.85\std{0.02}$ & $0.28\std{0.01}$ & $0.15\std{0.02}$ & $0.01\std{0.00}$ \\

\bottomrule
\end{tabular}
}


\resizebox{\columnwidth}{!}{
\begin{tabular}{crrrrrrrrrrrr}
\toprule
& \multicolumn{4}{c}{FC-20} & \multicolumn{4}{c}{FE-10} & \multicolumn{4}{c}{FE-20} \\ 
\cmidrule(lr){2-5} \cmidrule(lr){6-9} \cmidrule(lr){10-13}
$t$ & 
\multicolumn{1}{c}{F-each} & \multicolumn{1}{c}{G-each} & \multicolumn{1}{c}{F-total} & \multicolumn{1}{c}{G-total} & 
\multicolumn{1}{c}{F-each} & \multicolumn{1}{c}{G-each} & \multicolumn{1}{c}{F-total} & \multicolumn{1}{c}{G-total} &
\multicolumn{1}{c}{F-each} & \multicolumn{1}{c}{G-each} & \multicolumn{1}{c}{F-total} & \multicolumn{1}{c}{G-total} \\
\midrule

$2$ &
% FC-20
$0.41\std{0.07}$ & $0.10\std{0.01}$ & $0.29\std{0.04}$ & $0.10\std{0.01}$ &
% FE-10
$0.51\std{0.15}$ & $0.08\std{0.02}$ & $0.34\std{0.11}$ & $0.07\std{0.02}$ &
% FE-20
$0.49\std{0.08}$ & $0.09\std{0.01}$ & $0.33\std{0.08}$ & $0.08\std{0.01}$ \\

$T/2$ & 
% FC-20
$0.79\std{0.08}$ & $0.24\std{0.01}$ & $0.12\std{0.01}$ & $0.01\std{0.00}$ &
% FE-10
$0.72\std{0.09}$ & $0.17\std{0.04}$ & $0.25\std{0.07}$ & $0.02\std{0.00}$ &
% FE-20
$0.75\std{0.04}$ & $0.26\std{0.02}$ & $0.13\std{0.01}$ & $0.01\std{0.00}$ \\

$T$ &
% FC-20
$0.84\std{0.06}$ & $0.30\std{0.01}$ & $0.08\std{0.01}$ & $0.00\std{0.00}$ &
% FE-10
$0.86\std{0.02}$ & $0.24\std{0.05}$ & $0.11\std{0.01}$ & $0.01\std{0.00}$ &
% FE-20
$0.89\std{0.00}$ & $0.31\std{0.01}$ & $0.09\std{0.03}$ & $0.00\std{0.00}$ \\

\bottomrule
\end{tabular}
}


\end{table}






\begin{figure}[t]

\centering
\begin{minipage}{.3\hsize}
\centering
\includegraphics[width=\hsize]{images/repr_c10_t_leg.pdf}
\subcaption{Fine-tuning for TinyImageNet after CL for C-10}
\end{minipage}
\hfill
\begin{minipage}{.3\hsize}
\centering
\includegraphics[width=\hsize]{images/repr_i100_c.pdf}
\subcaption{Fine-tuning for CIFAR100 after CL for I-100}
\end{minipage}
\hfill
\begin{minipage}{.3\hsize}
\centering
\includegraphics[width=\hsize]{images/repr_t10_c.pdf}
\subcaption{Fine-tuning for CIFAR100 after CL for T-10}
\end{minipage}

\caption{
The learning of representation through continual learning.
The x-axis means the number of tasks learned in continual learning (CL).
The pair of a CL/non-CL for (a), (b) and (c) is C-10/TinyImageNet, I-100/CIFAR100, and T-10/CIFAR100, respectively.
}
\label{fig:cl_representation_app}

\end{figure}



\section{Representation Learning in Continual Learning} \label{sec:app_representation}






All results (i.e., the pairs of a CL and non-CL dataset are C-10/TinyImageNet, I-100/CIFAR100, and T-10/CIFAR100) are presented in \cref{fig:cl_representation_app}.
We can see that SPG learns better representations in continual learning than baselines in all cases.









\section{Network Size} \label{sec:app_network_size}


The number of learnable parameters of each system is presented in \cref{tab:app_network_size}. Note that all approaches adopt AlexNet as their backbone, and the number of parameters vary depending on their additional structures such as attention mechanisms or sub-modules. It also depends on datasets because each dataset has a different number of tasks and in TIL, each task has a different classification head and the number of units in each classification head depends on the number of classes in each task.
It can be seen that CAT and SupSup need more parameters than SPG and other approaches.


\begin{table}[ht]
\centering
\caption{
The number of learnable parameters of each model. 
``M'' means a million ($1,000,000$).
}
\label{tab:app_network_size}
\begin{tabular}{crrrrrrrrr}
\toprule
& \multicolumn{5}{c}{Dissimilar tasks} & \multicolumn{4}{c}{Similar tasks} \\
\cmidrule(lr){2-6} \cmidrule(lr){7-10}
Model & \multicolumn{1}{c}{C-10} & \multicolumn{1}{c}{C-20} & \multicolumn{1}{c}{T-10} & \multicolumn{1}{c}{T-20} & \multicolumn{1}{c}{I-100} &
\multicolumn{1}{c}{FC-10} & \multicolumn{1}{c}{FC-20} & \multicolumn{1}{c}{FE-10} & \multicolumn{1}{c}{FE-20} \\
\midrule

(MTL) & 
$6.7$M & $6.7$M & 
$6.9$M & $6.9$M &
$8.6$M & 
$6.5$M & $6.6$M &
$7.7$M & $9.0$M \\

(ONE) &
$6.7$M & $6.7$M & 
$6.9$M & $6.9$M &
$8.6$M & 
$6.5$M & $6.6$M &
$7.7$M & $9.0$M \\

\hdashline

NCL &
$6.7$M & $6.7$M & 
$6.9$M & $6.9$M &
$8.6$M & 
$6.5$M & $6.6$M &
$7.7$M & $9.0$M \\

A-GEM &
$6.7$M & $6.7$M & 
$6.9$M & $6.9$M &
$8.6$M & 
$6.5$M & $6.6$M &
$7.7$M & $9.0$M \\

PGN &
$6.7$M & $6.7$M &
$6.7$M & $6.7$M &
$8.3$M & 
$6.0$M & $6.6$M &
$7.5$M & $8.9$M \\

PathNet &
$6.6$M & $6.8$M & 
$6.7$M & $6.6$M &
$8.4$M & %
$6.4$M & $6.4$M &
$7.8$M & $8.7$M \\

HAT &
$6.8$M & $6.8$M &
$7.0$M & $7.0$M &
$9.0$M & 
$6.6$M & $6.7$M &
$7.8$M & $9.1$M \\

CAT &
$39.5$M & $39.7$M &
$40.8$M & $40.9$M &
N/A & 
$38.5$M & $38.9$M &
$46.2$M & $55.1$M \\

SupSup &
$65.2$M & $130.2$M &
$65.4$M & $130.4$M &
$652.0$M &
$65.0$M & $130.1$M &
$65.8$M & $131.7$M \\

UCL &
$6.7$M & $6.7$M &
$6.9$M & $6.9$M &
$8.6$M &
$6.5$M & $6.6$M &
$7.7$M & $9.0$M \\

SI &
$6.7$M & $6.7$M & 
$6.9$M & $6.9$M &
$8.6$M & 
$6.5$M & $6.6$M &
$7.7$M & $9.0$M \\

TAG &
$6.7$M & $6.7$M & 
$6.9$M & $6.9$M &
$8.6$M & 
$6.5$M & $6.6$M &
$7.7$M & $9.0$M \\

WSN &
$6.7$M & $6.7$M &
$6.8$M & $6.8$M &
$8.5$M & 
$6.5$M & $6.5$M &
$7.7$M & $8.9$M \\

EWC &
$6.7$M & $6.7$M & 
$6.9$M & $6.9$M &
$8.6$M & 
$6.5$M & $6.6$M &
$7.7$M & $9.0$M \\

\hdashline

EWC-GI &
$6.7$M & $6.7$M & 
$6.9$M & $6.9$M &
$8.6$M & 
$6.5$M & $6.6$M &
$7.7$M & $9.0$M \\

SPG-FI &
$6.7$M & $6.7$M & 
$6.9$M & $6.9$M &
$8.6$M & 
$6.5$M & $6.6$M &
$7.7$M & $9.0$M \\

\midrule

\textbf{SPG} &
$6.7$M & $6.7$M & 
$6.9$M & $6.9$M &
$8.6$M & 
$6.5$M & $6.6$M &
$7.7$M & $9.0$M \\

\bottomrule
\end{tabular}
%}
\end{table}


